%\documentclass[12pt,a4paper]{article}  % Use this line if this document will be released
\documentclass[12pt,a4paper,draft]{article}  % Use this line if this document is a draft
\usepackage{ifdraft}


%% Bibliography
\usepackage{etoolbox}
\newcommand{\bibfile}{\jobname.bib}  % Name of the BibTeX file.
% ref.bib should be a symbolic link to the universal BibTeX file, which should be a local copy of
% https://github.com/equipez/bibliographie/blob/main/ref.bib
% Run `getbib` in the current directory under the draft mode to get the BibTeX file containing only
% the cited references. The name will be xyz.bib if this TeX file is xyz.tex.
\newcommand{\universalbib}{ref.bib}
\ifdraft{\IfFileExists{\universalbib}{\renewcommand{\bibfile}{\universalbib}}{}}{}
% The counter `cite' is used to count the number of citations.
\newcounter{cite}
\pretocmd{\cite}{\stepcounter{cite}}{}{}


%% Add line numbers in draft mode
\RequirePackage[mathlines]{lineno}
\ifdraft{\linenumbers}{}
\renewcommand{\linenumberfont}{\normalfont\scriptsize\sffamily\color{gray}}
\setlength{\linenumbersep}{\marginparsep}


%% Geometry
%\voffset=-1.5cm \hoffset=-1.4cm \textwidth=16cm \textheight=22.0cm  % Luis' setting
\usepackage[a4paper, textwidth=16.0cm, textheight=22.0cm]{geometry}
\renewcommand{\baselinestretch}{1.2}


%% Basic packages
\usepackage{amsmath,amsthm,amssymb,amsfonts}
\usepackage{mathtools}  % Provides \coloneqq
\usepackage{empheq}
\usepackage{xcolor}
\usepackage[bbgreekl]{mathbbol}
\DeclareSymbolFontAlphabet{\mathbbm}{bbold}
\DeclareSymbolFontAlphabet{\mathbb}{AMSb}
\usepackage{bbm}
\usepackage{upgreek}
\usepackage{accents}
\usepackage{xspace}
\usepackage{rotating}
\usepackage{multirow,booktabs}
\usepackage[en-US]{datetime2}


%% Format of the table of content
\usepackage[normalem]{ulem}
\usepackage[toc,page]{appendix}
\renewcommand{\appendixpagename}{\Large{Appendix}}
\renewcommand{\appendixname}{Appendix}
\renewcommand{\appendixtocname}{Appendix}
%\usepackage{sectsty}
\setcounter{tocdepth}{2}


%% Section title style
\usepackage{sectsty}
\sectionfont{\large}
\subsectionfont{\large}


%% Some colors
\definecolor{darkblue}{rgb}{0,0.1,0.5}
\definecolor{darkgreen}{rgb}{0,0.5,0.1}
\definecolor{darkyellow}{rgb}{0.65,0.65,0.01}


%% Todo notes
\ifdraft{
    \setlength{\marginparwidth}{2.42cm}
    \usepackage[tickmarkheight=3pt,textsize=small,backgroundcolor=blue!16,linecolor=purple,bordercolor=purple]{todonotes}
}{
    \newcommand{\todo}[1]{}
    \newcommand{\listoftodos}{}
}


%% Graph, tikz and pgf
%\usepackage{subfigure}
\setlength{\unitlength}{1mm}
% The \unitlength command is a Length command. It defines the units used in the Picture Environment.
\usepackage{graphicx}
%\usepackage{tikz,tikzscale,pgf,pgfarrows,pgfnodes,filecontents,tikz-cd}
\usepackage{tikz,tikzscale,pgf}
\usetikzlibrary{arrows,arrows.meta,patterns,positioning,decorations.markings,shapes}
\usepackage{pgfplots}
\usepackage{pgfplotstable}
\usepackage[justification=centering]{caption}
\usepgfplotslibrary{fillbetween}
\pgfplotsset{compat=1.11}


%% Turn off some unharmful warnings in draft mode
%% N.B.: DO NOT use `silence` together with `hyperref`. They will cause an infinite loop.
\ifdraft{
    \usepackage{silence}
    \WarningFilter{xcolor}{Incompatible color definition on}
    \WarningFilter{hyperref}{Draft mode on}
    \WarningFilter{refcheck}{Unused label}
    \WarningFilter{microtype}{`draft' option active}
    \WarningFilter{latex}{Writing or overwriting file} % Mute the warning about 'writing/overwriting file'
    \WarningFilter{latex}{Writing file} % Mute the warning about 'writing/overwriting file'
    \WarningFilter{latex}{Tab has} % Mute the warning about 'Tab has been converted to Blank Space'
    \WarningFilter{latex}{Marginpar on page} % Mute the warning about 'Marginpar on page xx moved'
    \WarningFilter{latex}{author given} % Mute the warning about 'No \author given'
}{}


%% Hyperref, url, and email
%% N.B.: DO NOT use `silence` together with `hyperref`. They will cause an infinite loop.
\ifdraft{\usepackage{refcheck}\newcommand{\url}{\texttt}}{
    \usepackage{hyperref}
    \hypersetup{colorlinks, linkcolor=darkblue, anchorcolor=darkblue, citecolor=darkblue, urlcolor=darkblue}
    \usepackage{url}
} % Check unused labels
\newcommand{\email}{\texttt}


%% Enumerate and itemize
\usepackage{eqlist}
\usepackage{enumitem}
\setlist[itemize]{leftmargin=*}
\setlist[enumerate]{leftmargin=*,label=\normalfont{(\alph*)}}


%% Algorithm environment
\usepackage[section]{algorithm}
\usepackage{algpseudocode,algorithmicx}
\newcommand{\INPUT}{\textbf{Input}}
\newcommand{\FOR}{\textbf{For}~}
\algrenewcommand\algorithmicrequire{\textbf{Input:}}
\algrenewcommand\algorithmicensure{\textbf{Output:}}
\algrenewcommand\alglinenumber[1]{\normalsize #1.}
\newcommand*\Let[2]{\State #1 $=$ #2}


%% Theorem-like environments
\newtheorem{theorem}{Theorem}[section]
\newtheorem{conjecture}{Conjecture}[section]
\newtheorem{corollary}{Corollary}[section]
\newtheorem{exercise}{Exercise}[section]
\newtheorem{lemma}{Lemma}[section]
\newtheorem{problem}{Problem}[section]
\newtheorem{proposition}{Proposition}[section]
\newtheorem{assumption}{Assumption}[section]
\newtheorem{example}{Example}[section]
\newtheorem{question}{Question}[section]
% Change theoremstyle to ``definition'', which uses textnormal for the text.
\theoremstyle{definition}
\newtheorem{definition}{Definition}[section]
\newtheorem{remark}{Remark}[section]
% proof
\usepackage{xpatch}
\xpatchcmd{\proof}{\itshape}{\normalfont\proofnamefont}{}{}
\newcommand{\proofnamefont}{\bfseries}

%% Equation numbering
\numberwithin{equation}{section}


%% Fine tuning
\usepackage{microtype}
\usepackage[nobottomtitles*]{titlesec} % No section title at the bottom of pages
% Prevent footnote from running to the next page
\interfootnotelinepenalty=10000
% No line break in inline math
\interdisplaylinepenalty=10000
\relpenalty=10000
\binoppenalty=10000
% No widow or orphan lines
\clubpenalty=10000
\widowpenalty=10000
\displaywidowpenalty=10000


% Use @ to put 1 math unit (mu) in math
% See https://nhigham.com/2013/01/07/fine-tuning-spacing-in-latex-equations/
% and also TeXbook p. 155.
\mathcode`@="8000{\catcode`\@=\active\gdef@{\mkern1mu}}


%% Operators, commands
\usepackage{relsize}
\usepackage{nccmath}
%\DeclareMathOperator*{\mcap}{\,\medmath{\bigcap}\,}
%\DeclareMathOperator*{\mcup}{\,\medmath{\bigcup}\,}
\DeclareMathOperator*{\mcap}{\,\mathsmaller{\bigcap}\,}
\DeclareMathOperator*{\mcup}{\,\mathsmaller{\bigcup}\,}
%\renewcommand{\cap}{\mcap}
%\renewcommand{\cup}{\mcup}

\newcommand{\ceil}[1]{ {\lceil{#1}\rceil} }
\newcommand{\floor}[1]{ {\lfloor{#1}\rfloor} }

\DeclareMathOperator{\tr}{tr}
\DeclareMathOperator{\sort}{sort}
\DeclareMathOperator*{\Argmax}{Argmax}
\DeclareMathOperator*{\Argmin}{Argmin}
\DeclareMathOperator*{\Arglocmin}{Arglocmin}
\DeclareMathOperator*{\argmax}{argmax}
\DeclareMathOperator*{\argmin}{argmin}
\DeclareMathOperator*{\diag}{diag}
\DeclareMathOperator*{\Diag}{Diag}
\DeclareMathOperator{\Span}{span}
\DeclareMathOperator{\med}{med}
\DeclareMathOperator{\essinf}{essinf}
\DeclareMathOperator{\cl}{cl}
\DeclareMathOperator{\vol}{vol}
\DeclareMathOperator{\comp}{C}
\DeclareMathOperator{\sign}{sign}
\DeclareMathOperator{\rank}{rank}
\DeclareMathOperator{\range}{range}
\DeclareMathOperator{\card}{card}
\DeclareMathOperator{\diam}{diam}
\DeclareMathOperator{\dist}{dist}
\newcommand{\disth}{{\operatorname{\updelta_{\sss{H}}}}}
\newcommand{\ind}{\mathbbm{1}}
%\newcommand*{\defeq}{\stackrel{\mbox{\normalfont\tiny{\textnormal{def}}}}{=}}
\newcommand\defeq{\mathrel{\overset{\makebox[0pt]{\mbox{\normalfont\tiny\sffamily def}}}{=}}}

\newcommand{\RR}{\mathbb{R}}
\newcommand{\BB}{\mathcal{B}}
\renewcommand{\SS}{\mathbb{S}}
\newcommand{\TT}{\mathcal{T}}
\newcommand{\ZZ}{\mathbb{Z}}
\newcommand{\NN}{\mathbb{N}}
\newcommand{\FF}{\mathcal{F}}
\newcommand{\CC}{\mathbb{C}}
\newcommand{\XX}{\mathcal{X}}
\newcommand{\sset}{\mathcal{S}}
\newcommand{\pen}{h}
\newcommand{\penpar}{\mu}
\newcommand{\res}{\rho}
\newcommand{\col}{r}
\newcommand{\ofd}{\mathcal{F}}
\newcommand{\stf}[1]{\mathbb{S}^{#1}}
\newcommand{\sss}[1]{{\scriptscriptstyle{#1}}}
\newcommand{\sK}{{\scriptscriptstyle{K}}}
\newcommand{\sT}{{\scriptscriptstyle{T}}}
\newcommand{\fro}{{\scriptstyle{\textnormal{F}}}}
\newcommand{\trs}{{\scriptstyle{\mathsf{T}}}}
\newcommand{\hmt}{{\scriptstyle{{\mathsf{H}}}}}
\newcommand{\pin}{{\scriptstyle{{\mathsf{+}}}}}
\newcommand{\inv}{{-1}}
\newcommand{\adj}{*}
\newcommand{\ones}{\mathbf{1}}

\newcommand{\cs}{\text{c}}
\newcommand{\hp}{\circ}
\newcommand{\cc}{\sss{\textnormal{C}}}
\newcommand{\dec}{\sss{\textnormal{D}}}
\newcommand{\cauchy}{\sss{\textnormal{C}}}
\newcommand{\scauchy}{\sss{\textnormal{S}}}
\newcommand{\crit}{\textnormal{crit}}
\newcommand{\rsg}{\hat{\partial}}
\newcommand{\gsg}{\partial}
\newcommand{\dom}{\textnormal{dom}}
\newcommand{\tf}{{\textnormal{f}}}
\newcommand{\tg}{{\textnormal{g}}}
\newcommand{\ts}{{\textnormal{s}}}
\newcommand{\st}{\textnormal{s.t.}}
\newcommand{\etc}{{etc.}\xspace}
\newcommand{\ie}{{i.e.}\xspace}
\newcommand{\eg}{{e.g.}\xspace}
\newcommand{\etal}{{et al.}\xspace}
\newcommand{\iid}{\text{i.i.d.}\xspace}
\newcommand{\as}{\text{a.s.}\xspace}

\newcommand{\me}{\mathrm{e}}
\newcommand{\md}{\mathrm{d}}
\newcommand{\mi}{\mathrm{i}}
\newcommand{\lev}{\mathrm{lev}}
\newcommand{\bA}{\mathbf{A}}
\newcommand{\bx}{\mathbf{u}}
%\newcommand{\bb}{\mathbf{f}}
\newcommand{\bb}{\mathbf{r}}
\newcommand{\nov}{n_{\textnormal{o}}}
\xspaceaddexceptions{]\}}
% tex.stackexchange.com/questions/15252/why-does-xspace-behave-differently-for-parenthesis-vs-braces-brackets
\newcommand{\MATLAB}{\textsc{Matlab}\xspace}
\newcommand{\octave}{\mbox{GNU Octave}\xspace}
\newcommand{\prblm}{\texttt}
\DeclareMathAlphabet{\mathsfit}{T1}{\sfdefault}{\mddefault}{\sldefault}
\SetMathAlphabet{\mathsfit}{bold}{T1}{\sfdefault}{\bfdefault}{\sldefault}
\newcommand{\prbb}{\mathsfit{p}}
\newcommand{\pp}{\mathsf{p}}
\newcommand{\qq}{\mathsf{q}}
\newcommand{\ttt}{\mathsfit{t}}
\newcommand{\tol}{\varepsilon}
\newcommand{\bt}{\mathbf{t}}
\newcommand{\br}{\mathbf{r}}
\newcommand{\dd}{\mathbf{d}}
\newcommand{\ii}{\mathbf{i}}
\newcommand{\jj}{\mathbf{j}}
\newcommand{\xx}{\mathbf{x}}
\renewcommand{\pp}{\mathbf{p}}
\renewcommand{\ggg}{\mathbf{g}}
\newcommand{\GG}{\mathbf{G}}
\DeclareMathOperator{\expc}{\mathbb{E}}
\renewcommand{\Pr}{\mathbb{P}}
\newcommand{\lb}{\underline}
\newcommand{\ub}{\overline}

% mathlcal font
\DeclareFontFamily{U}{dutchcal}{\skewchar\font=45 }
\DeclareFontShape{U}{dutchcal}{m}{n}{<-> s*[1.0] dutchcal-r}{}
\DeclareFontShape{U}{dutchcal}{b}{n}{<-> s*[1.0] dutchcal-b}{}
\DeclareMathAlphabet{\mathlcal}{U}{dutchcal}{m}{n}
\SetMathAlphabet{\mathlcal}{bold}{U}{dutchcal}{b}{n}

% mathscr font (supporting lowercase letters)
%\usepackage[scr=dutchcal]{mathalfa}
%\usepackage[scr=esstix]{mathalfa}
%\usepackage[scr=boondox]{mathalfa}
%\usepackage[scr=boondoxo]{mathalfa}
\usepackage[scr=boondoxupr]{mathalfa}
%\newcommand{\model}{\mathscr{h}}
\newcommand{\model}{\tilde{f}}
\newcommand{\rmod}{F}

\newcommand{\Set}[1]{\mathcal{#1}}
\DeclareMathAlphabet{\mathpzc}{OT1}{pzc}{m}{it} % The mathpzc font
\newcommand{\slv}{\mathpzc}
% mathpzc looks great, but it stops working on 19 Feb 2020 for no reason.
%\newcommand{\slv}{\mathscr}
\newcommand{\software}{\texttt}
\DeclareMathOperator{\eff}{\mathsf{e}\;\!}
\DeclareMathOperator{\Eff}{\mathsf{E}\;\!}
\newcommand{\out}{{\text{out}}}


%% Commands for revision
\newcommand{\red}[1]{\textcolor{red}{#1}}
\newcommand{\blue}[1]{\textcolor{blue}{#1}}
\newcommand{\green}[1]{\textcolor{darkgreen}{#1}}
\newcommand{\TYPO}[1]{{\color{orange}{#1}}}
\newcommand{\MISTAKE}[1]{{\color{violet}{#1}}}
\newcommand{\REPHRASE}[1]{{\color{darkgreen}{#1}}}
\newcommand{\REVISE}[1]{{\color{blue}{#1}}}
\newcommand{\REVISEred}[1]{{\color{red}{#1}}}
\newcommand{\COMMENT}{\todo}  % Needs the todonotes package
%\newcommand{\COMMENT}[1]{\textcolor{brown}{{\small{(comment: #1)}}}}  % This puts comments inline

% Use the following if revision is finished
%\newcommand{\TYPO}{}
%\newcommand{\MISTAKE}{}
%\newcommand{\REPHRASE}{}
%\newcommand{\REVISE}{}
%\newcommand{\REVISEred}{}
%\newcommand{\COMMENT}[1]{}  % Input ignored.


%%%%%%%%%%%%%%%%%%%%%%%%%%%%%%%%%%%%%%%%%%%%%%%%%%%%%%%%%%%%%%%%%%%%%%%%%%%%%%%%%%%%%%%%%%%%%%%%%%%%
\title{Course Notes on Optimization}

\date{\DTMnow}

\author{
    Zhu Huatao
    \thanks{}
    %\and
    %Author2
    %\thanks{Information2}
}


\begin{document}

\maketitle

%\begin{abstract}
%\end{abstract}

%\textbf{Keywords}: Keyword1, Keyword2
%%%%%%%%%%%%%%%%%%%%%%%%%%%%%%%%%%%%%%%%%%%%%%%%%%%%%%%%%%%%%%%%%%%%%%%%%%%%%%%%%%%%%%%%%%%%%%%%%%%%

\section{Basic Properties}

\subsection{Convex Set}

Here we will talk about a kind of subset~$A$ of~$\RR^n$.

\begin{definition} \label{def1.1}
    For any~$x,y$ in~$A$ and any~$a \in \left[ 0,1 \right]$, if 
    \[
    a x + (1-a) y \in A,
    \]
    then we say that~$A$ is a \textbf{convex set}.
\end{definition}

This definition implies that for any~$x,y$ in a convex set~$A$, the line segment connecting~$x$ and~$y$
is in~$A$. Here are some example of convex sets.

\begin{enumerate}
    \item An ellipsoid is a convex set. An ellipsoid in~$\RR^n$ can be expressed by
    $$ x^T B x \leq 1, $$
    where~$B$ is positive definite and symmetric.

    \item A hyperplane can be expressed by
    $$ a^T x = b, $$
    where~$a$ and~$b$ are fixed vector.
    Any hyperplane and the two half-spaces separated by it are convex sets.

\end{enumerate}

It's straightforward by definition~\eqref{def1.1} that we have following facts of convex sets.

\begin{proposition}
    For any two convex sets~$A$ and~$B$, the intersection~$A \cap B$ is a convex set.
\end{proposition}

\begin{proposition}
    An affine transformation is the composition of a linear transformation and a translation.
    Every affine transformation~$T: \RR^n \rightarrow \RR^n$ can be expressed by
    $$ Tx=Ax+b ,$$
    where~$A \in \RR^{n \times n}$ and~$b \in \RR^n$. We have two following conclusions.
    \begin{enumerate}
        \item If~$A$ is a convex set, then~$TA \in \RR^{n \times n}$ is a convex set.
        \item If~$S$ is a convex set, then~$T^{-1}S \in \RR^{n \times n}$ is a convex set.
    \end{enumerate}
\end{proposition}

\begin{proposition}
    For any two convex sets~$A$ and~$B$, the sum~$A+B$ is a convex set. 
\end{proposition}

\begin{proposition}
    For any two convex sets~$A \in \RR^n$ and~$B \in \RR^m$, we have 
    that~$A \times B$ is convex.
\end{proposition}

For any three elements~$x_1,x_2$ and~$x_3$ of a convex set~$A$, we have
$$ a_1 x_1 + a_2 x_2 + a_3 x_3 \in A, $$
where~$a_1,a_2,a_3 \in \left[ 0,1 \right]$ and~$a_1+a_2+a_3=1$.
By induction, we can extend this formula to any finite subset of~$A$.

\begin{definition}
    For any~$x_1,x_2,\cdots,x_n \in A \subset \RR^n$ and~$a_1,a_2,\cdots,a_n \in \left[0,1\right]$ 
    satisfying~$\sum_{i=1}^{n}a_i=1$, we call~$ \sum_{i=1}^{n}a_i x_i $
    a \textbf{convex combination} of~$A$.
\end{definition}

A subset~$A$ is convex if and only if any convex combination of~$A$ is an element of~$A$.
For any subset~$B$ of~$\RR^n$, we want to find a convex set~$C$ that~$B \subset C$ and any
convex set that contains~$B$ also contains~$C$. In other words, we want to find the "smallest" 
convex set that contains~$B$.

\begin{definition}
    For any set~$A \subset \RR^n$, we say that the \textbf{convex bull}~$A$ is the set of all
    convex combination of~$A$. We denote it as~$conv(A)$.
\end{definition}

It's straightforward by the definition of the convex bull that any convex bull is convex 
and~$conv(A)$ contains~$A$.

\begin{theorem}
    For any subset~$A \subset \RR^n$, the convex bull of~$A$ contains~$A$ and 
    any convex set that contains~$A$ also contains~$conv(A)$.
\end{theorem}
A straightforward corollary is that~$A=conv(A)$ if and only if~$A$ is a convex set.

\subsection{Affine Set}

\begin{definition}
    For any~$x_1,x_2,\cdots,x_n \in A \subset \RR^n$ and~$a_1,a_2,\cdots,a_n \in \RR$ 
    satisfying~$\sum_{i=1}^{n}a_i=1$, we call~$ \sum_{i=1}^{n}a_i x_i $
    an \textbf{affine combination} of~$A$.
\end{definition}

An affine combination is a generalization of a convex combination.

\begin{definition} \label{def1.5}
    For any~$A \subset \RR^n$, if any affine combination of~$A$ is an element of~$A$,
    we say that~$A$ is an \textbf{affine set}.
\end{definition}
If~$A$ is an affine set, then for any two elements~$x,y \in A$, the line passing 
through~$x$ and~$y$ is in~$A$.
It's trivial that any convex set is an affine set. Here are some examples of affine set.
\begin{enumerate}
    \item Any line is an affine set.
    \item Any subspace and any hyperplane are both affine sets.
\end{enumerate}
Those examples can be straightforwardly verified by definition~\eqref{def1.5}.
Similar to the convex bull, we also have affine bull.

\begin{definition}
    For any set~$A \subset \RR^n$, we say that the \textbf{affine bull}~$A$ is the set of all
    affine combination of~$A$. We denote it as~$aff(A)$.
\end{definition}

We have that~$aff(A)$ is an affine set and is the "smallest" affine set that
contains~$A$. It's trivial that~$A=aff(A)$ if and only if~$A$ is an affine set.
For any set~$A \subset \RR^n$, we can notice that
$$ A \subset conv(A) \subset aff(A). $$


\subsection{Convex Function}
\begin{definition} \label{def1.7}
    For any function~$f: D \to \RR$ where~$D \subset \RR^n$ is a convex set, if for any~$x,y \in D$ and any~$a \in \left[ 0,1 \right]$, we have
    $$ f(ax+(1-a)y) \leq a f(x) + (1-a) f(y), $$
    then we say that~$f$ is a \textbf{convex function}. If the above inequality is strict for any~$x \neq y$ and any~$a \in (0,1)$, then we say that~$f$ is a \textbf{strictly convex function}.
    If there exists~$\mu > 0$ such that for any~$x,y \in D$ and any~$a \in \left[ 0,1 \right]$, we have
    $$ f(ax+(1-a)y) \leq a f(x) + (1-a) f(y) - \frac{\mu}{2} a(1-a) \|x-y\|^2, $$
    then we say that~$f$ is a~\textbf{$\mu$-strongly convex function}.
\end{definition}
For the case of~$n=1$, the definition of convex function can be interpreted as that the line segment connecting any two points on the curve of~$f$ is above the curve. 
A equivalent definition of~$\mu$-strongly convex function is that~$f(x) - \frac{\mu}{2} \|x\|^2$ is a convex function. 
This can be straightforwardly verified by Definition~\eqref{def1.7} and the properties of $2$-norm.
Here are some examples of convex functions.
\begin{enumerate}
    \item Any affine function is a convex function. An affine function can be expressed by
    $$ f(x) = a^T x + b, $$
    where~$a \in \RR^n$ and~$b \in \RR$ are fixed vector and scalar.
    \item For any quadratic function~$f(x) = \frac{1}{2} x^T A x + b^T x + c$, where~$A \in \RR^{n \times n}$ is a positive semidefinite matrix,~$b \in \RR^n$ and~$c \in \RR$ are fixed vector and scalar, we have that~$f$ is a convex function.
\end{enumerate}

In convex analysis, we are often interested in the following function extension:
\[
\tilde{f}(x) = \begin{cases}
    f(x), & x \in D, \\
    +\infty, & x \notin D.
\end{cases}
\]
The convexity of~$f$ is deeply related to the convexity of~$\tilde{f}(x)$. In fact we have the following result.
\begin{proposition} \label{prop1.5}
    For any function~$f: D \to \RR$ where~$D \subset \RR^n$ is a convex set, the function~$\tilde{f}(x)$ defined above is convex if and only if~$f(x)$ is convex.
\end{proposition}
\begin{proof}
    The "only if" part is trivial. For the "if" part, for any~$x,y \in D$ and any~$a \in \left[ 0,1 \right]$, we have
    $$ f(ax+(1-a)y) = \tilde{f}(ax+(1-a)y) \leq a \tilde{f}(x) + (1-a) \tilde{f}(y) = a f(x) + (1-a) f(y). $$
    If at least one of~$x$ and~$y$ is not in~$D$, then the above inequality also holds since the right hand side is~$+\infty$.
    Thus we have that~$\tilde{f}(x)$ is a convex function.
\end{proof}

For any subset~$A$ of~$\RR^n$, we can define the indicator function~$\delta_A(x)$ as
\[ \delta_A(x) = \begin{cases}
    0, & x \in A, \\
    +\infty, & x \notin A.
\end{cases} 
\]
Similar to Proposition~\eqref{prop1.5}, we can prove that~$\delta_A(x)$ is a convex function if and only if~$A$ is a convex set. In fact, we have the following result.
\begin{proposition}
    For any subset~$A$ of~$\RR^n$, the indicator function~$\delta_A(x)$ is a convex function if and only if~$A$ is a convex set.
\end{proposition}
\begin{proof}
    The "if" part is trivial. For the "only if" part, for any~$x,y \in A$ and any~$a \in \left[ 0,1 \right]$, we have
    $$ \delta_A(ax+(1-a)y) \leq a \delta_A(x) + (1-a) \delta_A(y) = 0. $$
    Thus we have that~$ax+(1-a)y \in A$. Therefore, we have that~$A$ is a convex set.
\end{proof}

The above discussion focuses on general convex functions. We are interested in 
some convex funtions that are first-order or second-order continously differentiable.
In the following, we will give some equivalent characterizations of convex functions in terms of their first-order or second-order derivatives.
These characterizations are very useful in optimization since they can be used to verify the convexity of a function in a more convenient way.

We first give the first-order characterization of convex functions. The following result is also known as the first-order condition for convexity.
\begin{theorem}
    For any function~$f: D \to \RR$ where~$D \subset \RR^n$ is a convex set and~$f$ is continously differentiable, the following statements are equivalent.
    \begin{enumerate}
        \item The function~$f$ is a convex function.
        \item For any~$x,y \in D$, we have
        $$ f(y) \geq f(x) + \nabla f(x)^T (y-x). $$
        \item For any~$x,y \in D$, we have
        $$ (\nabla f(y) - \nabla f(x))^T (y-x) \geq 0. $$
    \end{enumerate}
\end{theorem}
\begin{proof}
    (a) $\Rightarrow$ (b): For any~$x,y \in D$ and any~$a \in \left[ 0,1 \right]$, we have
    \begin{align*}
        f(ax+(1-a)y) & = f(y+a(x-y)) \\
        & = f(y)+a \nabla f(y)^T (x-y) + o(\|x-y\|). \\
    \end{align*}
    From the above inequality and the convexity of~$f$, we have
    \[
    f(y)+a \nabla f(y)^T (x-y) + o(\|x-y\|) \leq a f(x) + (1-a) f(y).
    \]
    Thus we have 
    \[ f(y) \geq f(x) + \nabla f(y)^T (y-x). \]
    (b) $\Rightarrow$ (a):  For any~$x,y \in D$ and any~$a \in \left[ 0,1 \right]$, we denote~$z=ax+(1-a)y$. 
    We have
    \[
    f(x) \geq f(z) + \nabla f(z)^T (x-z)
    \]
    and 
    \[
    f(y) \geq f(z) + \nabla f(z)^T (y-z).
    \]
    Combining with the definition of~$z$, we have
    \begin{align*}
        a f(x) +(1-a) f(y) & \geq f(z) + \nabla f(z)^T (a(x-z)+(1-a)(y-z)) \\
        & = f(z) + \nabla f(z)^T (z-z)= f(z).
    \end{align*}
    (b) $\Rightarrow$ (c): For any~$x,y \in D$, we have
    \[
    f(x) \geq f(y) + \nabla f(y)^T (x-y)
    \]
    and 
    \[
    f(y) \geq f(x) + \nabla f(x)^T (y-x).
    \]
    Adding the above two inequalities, we have
    \[
    (\nabla f(y) - \nabla f(x))^T (y-x) \geq 0.
    \]
    (c) $\Rightarrow$ (b): For any~$x,y \in D$, we have
    \begin{align*}
        f(y) & = f(x) + \int_{0}^{1} \nabla f(x+t(y-x))^T (y-x) dt \\
        & = f(x) + \int_{0}^{1} \nabla f(x)^T (y-x) dt + \int_{0}^{1} (\nabla f(x+t(y-x)) - \nabla f(x))^T (y-x) dt \\
        & \geq f(x) + \nabla f(x)^T (y-x).
    \end{align*}
\end{proof}
The above result can be interpreted as that the tangent hyperplane of the curve of~$f$ at any point is a global underestimator of~$f$.
The following result is also known as the second-order condition for convexity.
\begin{theorem}
    For any function~$f: D \to \RR$ where~$D \subset \RR^n$ is a convex set and~$f$ is twice continously differentiable.
    Then~$f$ is convex if and only if the Hessian matrix~$\nabla^2 f(x)$ is positive semidefinite for any~$x \in D$.
\end{theorem}
\begin{proof}
    For the "if" part, for any~$x,y \in D$ and any~$a \in \left[ 0,1 \right]$, we have
    \[
    \nabla f(y) - \nabla f(x) = \int_{0}^{1} \nabla^2 f(x+t(y-x)) (y-x) dt.
    \]
    Since the Hessian matrix~$\nabla^2 f(x)$ is positive semidefinite for any~$x \in D$, we have
    \[
    (\nabla f(y) - \nabla f(x))^T (y-x) = \int_{0}^{1} (y-x)^T \nabla^2 f(x+t(y-x)) (y-x) dt \geq 0.
    \]
    This completes the proof of the "if" part.

    For the "only if" part, for any~$x,y \in D$ and any~$a \in \left[ 0,1 \right]$, we denote~$z=ax+(1-a)y$.
    We have
    \[ \nabla f(z) - \nabla f(x) = \int_{0}^{1} \nabla^2 f(x+t(z-x)) (z-x) dt. \]
    Since~$f$ is a convex function, we have
    \[ (\nabla f(z) - \nabla f(x))^T (z-x) = \int_{0}^{1} (z-x)^T \nabla^2 f(x+t(z-x)) (z-x) dt \geq 0. \]
    Thus we have
    \[ (z-x)^T \nabla^2 f(x+t(z-x)) (z-x) \geq 0 \]
    for any~$t \in [0,1]$. Since~$x$ and~$y$ are arbitrary, we have
    \[ (y-x)^T \nabla^2 f(x+t(y-x)) (y-x) \geq 0 \]
    for any~$x,y \in D$ and any~$t \in [0,1]$. This completes the proof of the "only if" part.
\end{proof}
Similarly, the second-order condition for~$\mu$-strong convexity is that~$\nabla^2 f(x) \succeq \mu I$ for any~$x \in D$.
We can also prove that~$\nabla^2 f(x) \succ  0$ implies the strict convexity of~$f$. However, the converse is not true. For example, the function~$f(x) = x^4$ is a strictly convex function but its second-order derivative is not positive definite at~$x=0$.

The composition of a convex function with certain operations is still convex.
\begin{proposition}
    For convex functions~$f,f_1,f_2,\cdots,f_n$, the following functions are all convex.
    \begin{enumerate}
        \item For any~$a_1,a_2,\cdots,a_n \geq 0$, the function~$g(x) = a_1 f_1(x) + a_2 f_2(x) + \cdots + a_n f_n(x)$ is a convex function.
        \item For any~$A \in \RR^{m \times n}$ and~$b \in \RR^m$, the function~$h(x) = f(Ax+b)$ is a convex function.
        \item The function~$k(x) = \max\{f_1(x),f_2(x)\}$ is a convex function.
    \end{enumerate}
\end{proposition}
\begin{proof}
    \begin{enumerate}
        \item It's straightforward by the definition of convex functions that~$g(x)$ is a convex function. We only need to notice that
        the intersection of finite number of convex sets is a convex set and thus the domain of~$g(x)$ is a convex set.
        \item Since the domain~$D$ of~$f$ is a convex set and the inverse image of a convex set under an affine mapping is convex, we have that the domain of~$h(x)$ is a convex set.
        For any~$x,y$ in the domain of~$h$ and any~$a \in \left[ 0,1 \right]$, we have
        \begin{align*}
            h(ax+(1-a)y) & = f(A(ax+(1-a)y)+b) \\
            & = f(a(Ax+b)+(1-a)(Ay+b)) \\
        \end{align*}
        From the convexity of~$f$, we have
        \begin{align*}
            h(ax+(1-a)y)   & \leq a f(Ax+b) + (1-a) f(Ay+b) \\
            & = a h(x) + (1-a) h(y).
        \end{align*}
        Thus we have that~$h(x)$ is a convex function.
        \item The same as above, we have that the domain of~$k(x)$ is a convex set. For any~$x,y$ in the domain of~$k$ and any~$a \in \left[ 0,1 \right]$, we have
        \begin{align*}
            a k(x) + (1-a) k(y) &= a \max\{f_1(x),f_2(x)\} + (1-a) \max\{f_1(y),f_2(y)\} \\
            & \geq a f_{i}(x) + (1-a) f_{i}(y). \\
        \end{align*}
        for~$i=1,2$.
        From the convexity of~$f_i$, we have
        \[
            a k(x) + (1-a) k(y) \geq f_{i}(ax+(1-a)y)
        \]
        for~$i=1,2$. Thus we have
        \[ a k(x) + (1-a) k(y) \geq \max\{f_1(ax+(1-a)y),f_2(ax+(1-a)y)\} = k(ax+(1-a)y). \]
        Thus we have that~$k(x)$ is a convex function.
    \end{enumerate}
\end{proof}




%%%%%%%%%%%%%%%%%%%%%%%%%%%%%%%%%%%%%%%%%%%%%%%%%%%%%%%%%%%%%%%%%%%%%%%%%%%%%%%%%%%%%%%%%%%%%%%%%%%%
%% References
% Include references only if there are citations.
\ifnum\value{cite}>0
    \small
    \bibliography{\bibfile}
    \bibliographystyle{plain}
\fi

%% The end
\end{document}
